% ---------------------------------------------------------------------------
%  Copyright / colophon
% ---------------------------------------------------------------------------
\thispagestyle{empty}
\vspace*{\fill}
\noindent
{\small
\textit{Playing with RISC-V: Assembly, C, and Retro Games on the PRG32 Platform.}\\
A first-year course companion for the Computer Architecture (assembly language)
and Computer Programming (C language) curricula.\par
\vspace{1em}

\noindent
This textbook is built around the \prg open educational platform, an
educational runtime for \riscv assembly and C games that targets the
ESP32\nobreakdash-C6 RISC\nobreakdash-V microcontroller for physical hardware
and an ESP32\nobreakdash-C3 Espressif QEMU target for desktop emulation. The
platform, its source code, and its documentation are distributed under the MIT
License by its authors.\par
\vspace{1em}

\noindent
The \prg software is the work of Raffaele Montella (academic supervisor and
project lead) and Ivan Cafiero (student contributor) at the University of Naples
``Parthenope''. Readers who use \prg in coursework, reports, theses, or teaching
material are asked to cite the project using the \fname{CITATION.cff} metadata
shipped in the repository \citep{prg32repo}.\par
\vspace{1em}

\noindent
RISC\nobreakdash-V is a registered trademark of RISC\nobreakdash-V
International. ESP32, ESP32\nobreakdash-C3, ESP32\nobreakdash-C6, and ESP\nobreakdash-IDF are
products and trademarks of Espressif Systems. All other trademarks are the
property of their respective owners and are used here for identification only.\par
\vspace{1em}

\noindent
Edition compiled \today. Typeset in \LaTeX{} with Computer Modern.\par
}
\vspace*{\fill}
\clearpage
